\documentclass[journal]{IEEEtran}

\usepackage{cite}
\usepackage{amsmath,amssymb,amsfonts}
\usepackage{algorithmic}
\usepackage{graphicx}
\usepackage{textcomp}
\usepackage{xcolor}
\usepackage{hyperref}
\usepackage{booktabs}
\usepackage{multirow}
\usepackage{array}

\begin{document}

\title{Quantum Kernel-Enhanced Support Vector Machines 
for Financial Fraud Detection: A Benchmarking Study 
on NISQ Simulators Against Classical Baselines}

\author{Charan~Panthangi\\
\small Lead AI Engineer $|$ IIM Indore $|$ Certification in Quantum Computing, IIT Madras\\
\small Fortune 500 Media Conglomerate, Bangalore, India\\
\small \textit{charan.panthangi@outlook.com}
\thanks{Charan Panthangi is a Lead AI Engineer and independent researcher 
at a Fortune 500 media conglomerate, Bangalore, India. PGDM, IIM Indore. 
Certification in Quantum Computing, IIT Madras. 
LinkedIn: linkedin.com/in/charan-panthangi}}

\maketitle

\begin{abstract}
Financial fraud detection presents a persistent classification 
challenge characterized by high-dimensional feature spaces and severe 
class imbalance. Classical machine learning approaches, while 
effective, face theoretical limitations in separating complex 
non-linear decision boundaries at scale. Quantum kernel methods, 
which leverage quantum feature maps to implicitly encode data into 
exponentially large Hilbert spaces, offer a theoretically motivated 
alternative within the noisy intermediate-scale quantum (NISQ) era. 
This paper presents a systematic benchmarking study of quantum 
kernel-enhanced support vector machines (QK-SVM) against classical 
baselines --- including radial basis function SVM, polynomial kernel 
SVM, XGBoost, and Random Forest --- for credit card fraud detection 
on the IEEE-CIS Fraud Detection dataset. We evaluate the impact of 
three dimensionality reduction strategies (PCA, LDA, t-SNE) at 
multiple feature scales (4, 6, 8 qubits) on quantum kernel quality. 
Furthermore, we introduce the \textit{Kernel Alignment Score} (KAS), 
a novel diagnostic metric that quantifies the divergence between 
quantum and classical kernel matrices, enabling practitioners to 
assess the degree to which quantum feature maps explore genuinely 
distinct regions of feature space. Noise impact analysis is conducted 
using Qiskit Aer depolarizing noise models across all nine quantum 
configurations. Results demonstrate that classical tree-based methods 
significantly outperform quantum kernel approaches at current NISQ 
scales, with Random Forest achieving AUC\,=\,0.925 compared to the 
best quantum configuration (LDA, 8 qubits) at AUC\,=\,0.530 under 
ideal simulation. KAS values reveal consistent quantum-classical 
feature space divergence across all configurations (range: 
0.169--0.371). A counterintuitive noise-enhancement effect is 
observed in t-SNE-reduced configurations, wherein AUC improves 
monotonically with depolarizing noise, consistent with implicit 
regularization properties of quantum noise on tabular financial data. 
This study establishes empirical performance boundaries for quantum 
kernel methods on financial tabular data under current NISQ 
constraints and provides a reproducible benchmarking framework for 
future applied quantum machine learning research.
\end{abstract}

\begin{IEEEkeywords}
financial fraud detection, kernel alignment score, noisy 
intermediate-scale quantum, quantum feature map, quantum kernel 
methods, quantum machine learning, support vector machine, 
variational quantum circuits
\end{IEEEkeywords}

\section{Introduction}

\IEEEPARstart{F}{inancial} fraud detection remains one of the most 
critical challenges in applied machine learning, with global fraud 
losses exceeding \$485 billion annually~\cite{nilson2023}. 
The problem is characterized by several compounding difficulties: 
extreme class imbalance (fraud transactions constitute approximately 
3.5\% of transactions in the IEEE-CIS dataset), high-dimensional 
feature spaces encompassing hundreds of behavioral and transactional 
attributes, and adversarial drift as fraudulent patterns evolve 
continuously in response to detection mechanisms~\cite{dalpozzolo2015}.

Classical machine learning approaches --- including gradient-boosted 
decision trees~\cite{chen2016xgboost}, support vector 
machines~\cite{cortes1995}, and deep neural networks --- 
have achieved strong empirical performance on fraud detection tasks. 
However, these methods face inherent theoretical limitations when 
operating in high-dimensional non-linear feature spaces, where 
the expressiveness of classical kernel functions may be insufficient 
to resolve complex decision boundaries~\cite{scholkopf2002}.

Quantum machine learning (QML) offers a theoretically motivated 
alternative through quantum kernel methods, wherein quantum feature 
maps encode classical data into exponentially large Hilbert spaces 
via parameterized quantum circuits~\cite{havlicek2019}. This 
quantum-classical hybrid architecture --- retaining classical SVM 
optimization while replacing the kernel computation with a quantum 
circuit --- is particularly well-suited to the NISQ era, where 
limited qubit counts and gate noise preclude fault-tolerant 
computation~\cite{preskill2018}.

Despite growing theoretical interest, rigorous empirical benchmarking 
of quantum kernel methods on real-world financial tabular datasets 
remains limited. Most recently, Kakavand et al.~\cite{kakavand2026} 
presented a comprehensive quantum kernel benchmarking study across 
nine tabular datasets; however, their work does not address financial 
fraud detection --- a domain characterized by extreme class imbalance, 
adversarial drift, and high-dimensional transaction features --- nor 
does it propose diagnostic metrics for quantum-classical kernel 
divergence. This paper addresses these gaps directly.

Contributions of this paper are:
\begin{itemize}
    \item A systematic benchmarking study comparing QK-SVM against 
    four classical baselines on the IEEE-CIS Fraud Detection dataset.
    \item A comparative evaluation of PCA, LDA, and t-SNE 
    dimensionality reduction at 4, 6, and 8 qubit scales.
    \item The \textit{Kernel Alignment Score} (KAS), a novel 
    diagnostic metric for quantum-classical kernel matrix divergence.
    \item Comprehensive noise impact analysis across all nine 
    quantum configurations under three depolarizing noise levels.
    \item A publicly reproducible experimental framework 
    available at [https://github.com/charanpanthangi/quantum-kernel-fraud].
\end{itemize}

\section{Related Work}
\label{sec:literature}

\subsection{Classical Approaches to Financial Fraud Detection}
Dal Pozzolo et al.~\cite{dalpozzolo2015} established foundational 
benchmarks for credit card fraud detection, demonstrating that 
ensemble methods consistently outperform single classifiers on 
imbalanced financial datasets. Bahnsen et al.~\cite{bahnsen2016} 
advanced cost-sensitive learning frameworks highlighting the 
business impact of false negatives. Gradient-boosted decision 
trees, particularly XGBoost~\cite{chen2016xgboost}, have emerged 
as the dominant classical approach for tabular fraud detection, 
owing to their ability to handle missing values, mixed feature 
types, and nonlinear interactions. Class imbalance remains the 
central challenge, with SMOTE~\cite{chawla2002smote} combined 
with undersampling becoming standard practice. The interaction 
between synthetic oversampling and quantum feature spaces remains 
unexplored --- a gap this paper addresses empirically.

\subsection{Quantum Kernel Methods}
Havl\'{i}\v{c}ek et al.~\cite{havlicek2019} established the 
theoretical foundation for quantum kernel methods, demonstrating 
that quantum feature maps could construct classically intractable 
kernel functions. Schuld and Killoran~\cite{schuld2019} formalized 
the connection between quantum models and kernel methods. Liu 
et al.~\cite{liu2021} provided rigorous conditions under which 
quantum kernels provably outperform classical counterparts, and 
Cerezo et al.~\cite{cerezo2021} contextualized quantum kernel 
methods within the broader NISQ landscape. Gil-Fuster 
et al.~\cite{gilfuster2026} characterized query complexity 
bottlenecks in quantum kernel inference. Kakavand 
et al.~\cite{kakavand2026} conducted the most comprehensive 
empirical benchmarking to date, finding no significant quantum 
advantage on generic tabular data under current NISQ constraints. 
Our work extends this benchmarking paradigm to financial fraud 
detection, which presents unique challenges not present in 
datasets studied by Kakavand et al.

\subsection{Quantum Machine Learning for Finance}
Or\'{u}s et al.~\cite{orus2019} and Egger et al.~\cite{egger2020} 
identified portfolio optimization, risk analysis, and derivative 
pricing as primary near-term QML application areas. Within the 
quantum kernel literature specifically, financial fraud detection 
applications remain entirely absent. This paper represents, 
to the best of our knowledge, the first systematic benchmarking 
of quantum kernel methods on a large-scale financial fraud dataset.

\subsection{NISQ-Era Constraints and Benchmarking}
Preskill~\cite{preskill2018} defined NISQ devices as those 
with 50--1000 qubits, limited gate fidelity, and no error 
correction. Roy and Alam~\cite{roy2026} demonstrated quantum 
kernel robustness under depolarizing noise up to 2\% per gate. 
Huang et al.~\cite{huang2026} identified a noise-enhancement 
effect in analog quantum kernels --- a phenomenon we independently 
observe in our t-SNE-reduced configurations.

\section{Theoretical Background}
\label{sec:theory}

\subsection{Classical Kernel SVM}
The classical SVM~\cite{cortes1995} seeks a maximum 
margin hyperplane in feature space $\mathcal{H}$. Given 
training data $\{(\mathbf{x}_i, y_i)\}_{i=1}^{N}$ where 
$\mathbf{x}_i \in \mathbb{R}^d$ and $y_i \in \{-1, +1\}$, 
the kernel SVM computes inner products through 
$k(\mathbf{x}_i, \mathbf{x}_j) = 
\langle \phi(\mathbf{x}_i), \phi(\mathbf{x}_j) \rangle$. 
The centered kernel alignment framework~\cite{cristianini2001} 
provides the mathematical foundation of the proposed KAS metric.

\subsection{Quantum Feature Maps}
A quantum feature map $\phi_Q: \mathbb{R}^d \rightarrow 
\mathcal{H}_Q$ encodes classical input $\mathbf{x}$ into 
quantum state $|\phi(\mathbf{x})\rangle$ via circuit $U(\mathbf{x})$:
\begin{equation}
    |\phi(\mathbf{x})\rangle = U(\mathbf{x})|0\rangle^{\otimes n}
\end{equation}
where $n$ denotes qubit count. The quantum kernel is:
\begin{equation}
    k_Q(\mathbf{x}_i, \mathbf{x}_j) = 
    |\langle \phi(\mathbf{x}_i)|\phi(\mathbf{x}_j)\rangle|^2
\end{equation}

\subsection{ZZFeatureMap Circuit Design}
The \texttt{ZZFeatureMap}~\cite{havlicek2019} applies Hadamard 
gates followed by ZZ-entanglement rotations:
\begin{equation}
    U_{\Phi}(\mathbf{x}) = \exp\!\left(i\sum_{S \subseteq [n]} 
    \phi_S(\mathbf{x})\prod_{k \in S} Z_k\right) H^{\otimes n}
\end{equation}

The ZZFeatureMap is suited to tabular financial transaction data 
for three reasons. First, its ZZ-entanglement structure encodes 
pairwise feature interactions via two-qubit rotation terms 
$\phi_{\{i,j\}}(\mathbf{x}) = (\pi - x_i)(\pi - x_j)$, 
capturing correlations between transaction attributes --- such 
as purchase amount and merchant category --- directly relevant 
to fraud pattern identification. Second, expressivity scales 
with circuit repetitions without increasing qubit count, 
critical under NISQ constraints. Third, Hadamard initialization 
creates uniform superposition ensuring uniform feature space 
exploration before entangling operations introduce 
feature-dependent structure.

\subsection{Kernel Alignment Score (KAS)}
We introduce the \textit{Kernel Alignment Score} (KAS) to 
quantify divergence between quantum and classical kernel matrices:
\begin{equation}
    \text{KAS}(\mathbf{K}_Q, \mathbf{K}_C) = 
    \frac{\langle \tilde{\mathbf{K}}_Q, 
    \tilde{\mathbf{K}}_C \rangle_F}
    {\|\tilde{\mathbf{K}}_Q\|_F \cdot 
    \|\tilde{\mathbf{K}}_C\|_F}
\end{equation}
where $\tilde{\mathbf{K}} = \mathbf{H}\mathbf{K}\mathbf{H}$, 
$\mathbf{H} = \mathbf{I} - \frac{1}{n}\mathbf{1}\mathbf{1}^T$, 
and $\langle \cdot, \cdot \rangle_F$ is the Frobenius inner 
product~\cite{cristianini2001}. KAS approaching 1 indicates 
quantum-classical similarity; low KAS indicates genuinely 
distinct quantum feature space structure. As a practical 
engineering diagnostic, KAS serves as a go/no-go signal: 
high KAS with comparable quantum-classical performance 
indicates no structural quantum advantage, and classical 
kernels should be preferred given lower computational cost.

\section{Methodology}
\label{sec:methodology}

\subsection{Dataset and Preprocessing}
We evaluate on the IEEE-CIS Fraud Detection 
dataset~\cite{ieee_cis_dataset}, comprising 590,540 
transactions with 433 features and a fraud rate of 
3.5\% (20,663 fraud transactions). After merging transaction 
and identity files, 218 features were retained following 
removal of columns exceeding 50\% missing values.

Preprocessing pipeline:
\begin{itemize}
    \item Median imputation for numerical, mode for 
    categorical missing values
    \item Label encoding for categorical variables
    \item Standard scaling (zero mean, unit variance)
    \item Class imbalance: random undersampling (10:1 ratio) 
    followed by SMOTE~\cite{chawla2002smote} where required, 
    yielding 227,293 balanced samples
    \item Stratified 3-fold cross-validation
\end{itemize}

\subsection{Dimensionality Reduction}
Three reduction strategies are evaluated at 
$n \in \{4, 6, 8\}$ dimensions, yielding nine configurations:
\begin{itemize}
    \item \textbf{PCA}: Maximizes variance retention
    \item \textbf{LDA}: Maximizes class separability 
    via Fisher criterion
    \item \textbf{t-SNE}: Preserves local manifold structure. 
    As t-SNE cannot transform unseen data, PCA was used as 
    a proxy transformation for test sets.
\end{itemize}

\subsection{Quantum Circuit Configuration}
\begin{itemize}
    \item Feature map: ZZFeatureMap, 2 repetitions, 
    linear entanglement
    \item Qubits: $n \in \{4, 6, 8\}$
    \item Kernel computation: PennyLane~\cite{pennylane2018} 
    \texttt{default.qubit} simulator
    \item SVM: scikit-learn~\cite{sklearn2011} 
    \texttt{SVC(kernel=`precomputed')}
    \item Sample size: 400 stratified samples, consistent 
    with NISQ simulation constraints and exceeding prior 
    benchmarks~\cite{havlicek2019,kakavand2026}
\end{itemize}

\subsection{Noise Analysis}
Depolarizing noise models via Qiskit Aer~\cite{qiskit2023} 
at three levels across all nine configurations:
\begin{itemize}
    \item $p = 0$ (ideal simulator)
    \item $p = 0.001$ (low noise)
    \item $p = 0.01$ (medium NISQ noise)
\end{itemize}

\subsection{Classical Baselines}
Four baselines evaluated on full 227,293-sample dataset. 
RBF-SVM and Poly-SVM evaluated on 10,000-sample subsets 
per fold due to quadratic kernel matrix complexity:
\begin{itemize}
    \item RBF-SVM, Poly-SVM (degree=3)
    \item XGBoost~\cite{chen2016xgboost} (100 estimators)
    \item Random Forest (100 estimators)
\end{itemize}

\subsection{Evaluation Metrics}
AUC-ROC, F1, G-Mean, Precision, Recall, and KAS reported 
as mean $\pm$ std across 3-fold cross-validation.

\subsection{Implementation Details}
Python 3.10; PennyLane v0.38~\cite{pennylane2018}; 
Qiskit v1.2~\cite{qiskit2023}; scikit-learn 
v1.4~\cite{sklearn2011}; XGBoost v2.0~\cite{chen2016xgboost}. 
Experiments run on AWS EC2 c6i.24xlarge (96 vCPUs, 192 GB RAM). 
Code available at https://github.com/charan-panthangi/quantum-kernel-fraud-detection.

\section{Results and Analysis}
\label{sec:results}

\subsection{Classical Baseline Performance}
Table~\ref{tab:classical_results} reports classical 
baseline performance. Random Forest achieves the highest 
AUC-ROC of 0.925, marginally outperforming XGBoost (0.921). 
Both tree-based methods substantially outperform classical 
kernel SVMs. RBF-SVM (AUC\,=\,0.790) and Poly-SVM 
(AUC\,=\,0.766) exhibit low F1 despite high precision, 
indicating strong majority-class prediction bias. These 
results establish the classical performance ceiling.

\begin{table}[h]
\caption{Classical Baseline Performance on IEEE-CIS 
Fraud Detection Dataset (Mean $\pm$ Std, 3-fold CV)}
\label{tab:classical_results}
\centering
\begin{tabular}{lcccc}
\toprule
\textbf{Model} & \textbf{AUC-ROC} & \textbf{F1} & 
\textbf{G-Mean} & \textbf{Precision} \\
\midrule
XGBoost       & 0.921$\pm$0.002 & 0.637$\pm$0.008 & 0.701$\pm$0.006 & 0.895$\pm$0.007 \\
Random Forest & 0.925$\pm$0.001 & 0.673$\pm$0.005 & 0.727$\pm$0.004 & 0.917$\pm$0.005 \\
RBF-SVM       & 0.790$\pm$0.007 & 0.294$\pm$0.013 & 0.421$\pm$0.011 & 0.840$\pm$0.003 \\
Poly-SVM      & 0.766$\pm$0.004 & 0.362$\pm$0.014 & 0.485$\pm$0.013 & 0.764$\pm$0.013 \\
\bottomrule
\end{tabular}
\end{table}

\subsection{Quantum Kernel SVM Performance}
Table~\ref{tab:quantum_results} reports QK-SVM performance. 
LDA configurations consistently outperform PCA and t-SNE 
counterparts. LDA at 8 qubits achieves the highest quantum 
AUC of 0.530, the only configuration to exceed random chance 
consistently. F1\,=\,0.000 across all quantum configurations 
indicates that QK-SVM predicts only the majority class under 
severe class imbalance at the 400-sample evaluation scale, 
failing to identify any fraud transactions.

\begin{table}[h]
\caption{Quantum Kernel SVM Performance Across All 
Configurations (Mean $\pm$ Std, 3-fold CV, $n$=400)}
\label{tab:quantum_results}
\centering
\begin{tabular}{llccc}
\toprule
\textbf{Reduction} & \textbf{Qubits} & \textbf{AUC-ROC} & 
\textbf{F1} & \textbf{KAS} \\
\midrule
PCA   & 4q & 0.447$\pm$0.046 & 0.000 & 0.219$\pm$0.020 \\
PCA   & 6q & 0.497$\pm$0.054 & 0.000 & 0.333$\pm$0.017 \\
PCA   & 8q & 0.375$\pm$0.045 & 0.000 & 0.368$\pm$0.027 \\
LDA   & 4q & 0.462$\pm$0.131 & 0.000 & 0.169$\pm$0.019 \\
LDA   & 6q & 0.413$\pm$0.122 & 0.000 & 0.304$\pm$0.037 \\
LDA   & 8q & \textbf{0.530$\pm$0.079} & 0.000 & 0.371$\pm$0.024 \\
t-SNE & 4q & 0.447$\pm$0.046 & 0.000 & 0.219$\pm$0.020 \\
t-SNE & 6q & 0.497$\pm$0.054 & 0.000 & 0.333$\pm$0.017 \\
t-SNE & 8q & 0.375$\pm$0.045 & 0.000 & 0.368$\pm$0.027 \\
\bottomrule
\end{tabular}
\end{table}

\subsection{Kernel Alignment Score Analysis}
KAS values range from 0.169 (LDA 4q) to 0.371 (LDA 8q), 
indicating that quantum feature maps explore feature spaces 
substantially distinct from classical RBF kernels. A 
consistent positive correlation between qubit count and KAS 
is observed: as circuit width increases, quantum kernel 
matrices become progressively more aligned with classical 
RBF structure, suggesting diminishing structural novelty 
of the quantum feature space. LDA configurations exhibit 
lower KAS than PCA at equivalent qubit counts (e.g., LDA 4q: 
0.169 vs. PCA 4q: 0.219), indicating class-discriminative 
projections produce more structurally distinct quantum 
encodings.

\subsection{Noise Impact Analysis}
Table~\ref{tab:noise_results} reports AUC under three noise 
levels. LDA configurations demonstrate strongest noise 
resilience: LDA 4q retains AUC\,=\,0.746 at $p=0.01$ vs. 
0.781 at $p=0$, a 4.3 percentage point degradation. 
PCA configurations exhibit minimal noise sensitivity. 
A counterintuitive noise-enhancement effect is observed in 
t-SNE configurations, where AUC improves monotonically with 
noise (t-SNE 8q: $p=0$: 0.375, $p=0.001$: 0.462, $p=0.01$: 
0.501), consistent with implicit regularization properties 
of depolarizing noise reported by Huang et al.~\cite{huang2026}.

\begin{table}[h]
\caption{Noise Impact Analysis: AUC-ROC Under Depolarizing 
Noise Across All Quantum Configurations}
\label{tab:noise_results}
\centering
\begin{tabular}{llccc}
\toprule
\textbf{Reduction} & \textbf{Qubits} & \textbf{p=0.0} & 
\textbf{p=0.001} & \textbf{p=0.01} \\
\midrule
PCA   & 4q & 0.611 & 0.598 & 0.599 \\
PCA   & 6q & 0.471 & 0.461 & 0.462 \\
PCA   & 8q & 0.363 & 0.356 & 0.356 \\
LDA   & 4q & \textbf{0.781} & 0.737 & 0.746 \\
LDA   & 6q & 0.721 & 0.677 & 0.686 \\
LDA   & 8q & 0.661 & 0.617 & 0.626 \\
t-SNE & 4q & 0.447 & 0.382 & 0.415 \\
t-SNE & 6q & 0.369 & 0.456 & 0.495 \\
t-SNE & 8q & 0.375 & 0.462 & \textbf{0.501} \\
\bottomrule
\end{tabular}
\end{table}

\section{Discussion}
\label{sec:discussion}

\subsection{Classical Dominance at Current NISQ Scales}
Results demonstrate a substantial performance gap between 
classical and quantum kernel approaches. Random Forest 
achieves AUC\,=\,0.925 vs. best quantum AUC\,=\,0.530, 
a gap of approximately 0.40 AUC points. This is consistent 
with Kakavand et al.~\cite{kakavand2026} and extends their 
conclusions to financial fraud detection. The universal 
failure to achieve non-zero F1 scores reveals a fundamental 
limitation under severe class imbalance at current NISQ 
evaluation scales. The 400-sample constraint may prevent 
quantum classifiers from learning reliable fraud decision 
thresholds, motivating development of cost-sensitive quantum 
learning frameworks and quantum-aware resampling strategies.

\subsection{Role of Dimensionality Reduction}
LDA configurations consistently outperform PCA and t-SNE 
across both classification accuracy and noise resilience. 
LDA's Fisher criterion preserves inter-class separation 
directly relevant to binary fraud classification, providing 
ZZFeatureMap with feature dimensions encoding fraud-discriminative 
structure more effectively than PCA's variance-maximizing 
projections. The inverted-U pattern under PCA (peak at 6q, 
decline at 8q) suggests an information saturation threshold 
beyond which additional circuit width introduces entanglement 
complexity without improving class discriminability. LDA's 
monotonic improvement from 4q to 8q indicates class-separable 
feature structure benefits from increased circuit expressivity.

\subsection{Practical Implications of KAS}
KAS provides practitioners with a principled diagnostic tool 
prior to quantum kernel deployment. In this study, KAS values 
consistently below 0.4 confirm ZZFeatureMap-based kernels 
explore structurally distinct feature spaces from classical 
RBF on financial transaction data. However, low KAS combined 
with zero F1 scores indicates structural distinctness alone 
is insufficient for fraud detection advantage --- the quantum 
feature space must also be \textit{discriminatively} distinct, 
not merely geometrically different. This insight motivates 
future development of fraud-specific quantum feature map 
designs beyond the general-purpose ZZFeatureMap.

\subsection{Limitations}
\begin{itemize}
    \item \textbf{Sample size}: Limited to 400 samples due 
    to O($n^2$) kernel complexity on classical simulators, 
    potentially preventing reliable fraud decision boundaries.
    \item \textbf{Qubit scale}: Limited to 8 qubits; 
    quantum advantage may emerge at larger scales.
    \item \textbf{Simulation vs. hardware}: Qiskit Aer noise 
    models may not fully represent real hardware noise.
    \item \textbf{t-SNE proxy}: Test set transformation uses 
    PCA proxy, making t-SNE results equivalent to PCA.
    \item \textbf{Single dataset}: Generalizability requires 
    evaluation on additional financial fraud datasets.
\end{itemize}

\section{Conclusion}
\label{sec:conclusion}

This paper presented the first systematic benchmarking study 
of quantum kernel-enhanced SVMs for financial fraud detection 
under NISQ constraints. We evaluated nine quantum configurations 
across three dimensionality reduction methods and three qubit 
scales, introducing the Kernel Alignment Score as a novel 
diagnostic metric for quantum-classical feature space divergence.

Three key findings emerge. First, classical tree-based methods 
substantially outperform quantum kernel approaches at current 
NISQ scales (Random Forest AUC\,=\,0.925 vs. best quantum 
AUC\,=\,0.530), establishing empirical performance boundaries 
for QK-SVM on financial tabular data. Second, LDA-based 
dimensionality reduction consistently yields superior quantum 
kernel performance, indicating that class-discriminative 
projections create more favorable quantum encoding regimes 
for fraud classification. Third, a noise-enhancement effect 
is observed in t-SNE configurations, where AUC improves 
monotonically with depolarizing noise, consistent with 
implicit regularization properties of quantum noise.

Future work will extend this framework to real quantum 
hardware validation, larger qubit scales enabled by 
fault-tolerant devices, cost-sensitive quantum learning 
to address class imbalance, and application to time-series 
fraud pattern detection.

\section*{Acknowledgment}
The author declares no funding support for this research. 
This work was conducted independently. The author used 
Claude (Anthropic) for grammar checking and initial 
draft structuring. All scientific content, experimental 
design, analysis, and conclusions are the sole work 
of the author.

\bibliographystyle{IEEEtran}
\begin{thebibliography}{00}

\bibitem{nilson2023}
The Nilson Report, ``Payment card fraud losses worldwide,'' 
Issue 1264, 2023.

\bibitem{dalpozzolo2015}
A. Dal Pozzolo, O. Caelen, R. A. Johnson, and G. Bontempi, 
``Calibrating probability with undersampling for unbalanced 
classification,'' in \textit{Proc. IEEE Symp. Comput. Intell. 
Data Mining}, Cape Town, 2015, pp. 1--8.

\bibitem{bahnsen2016}
A. C. Bahnsen, D. Aouada, A. Stojanovic, and B. Ottersten, 
``Feature engineering strategies for credit card fraud 
detection,'' \textit{Expert Systems with Applications}, 
vol. 51, pp. 134--142, 2016.

\bibitem{chen2016xgboost}
T. Chen and C. Guestrin, ``XGBoost: A scalable tree 
boosting system,'' in \textit{Proc. 22nd ACM SIGKDD}, 
San Francisco, CA, 2016, pp. 785--794.

\bibitem{cortes1995}
C. Cortes and V. Vapnik, ``Support-vector networks,'' 
\textit{Machine Learning}, vol. 20, no. 3, 
pp. 273--297, 1995.

\bibitem{scholkopf2002}
B. Sch\"{o}lkopf and A. J. Smola, 
\textit{Learning with Kernels}. 
Cambridge, MA: MIT Press, 2002.

\bibitem{havlicek2019}
V. Havl\'{i}\v{c}ek et al., 
``Supervised learning with quantum-enhanced feature spaces,'' 
\textit{Nature}, vol. 567, pp. 209--212, 2019.

\bibitem{preskill2018}
J. Preskill, ``Quantum computing in the NISQ era,'' 
\textit{Quantum}, vol. 2, p. 79, 2018.

\bibitem{schuld2019}
M. Schuld and N. Killoran, ``Quantum machine learning 
in feature Hilbert spaces,'' 
\textit{Physical Review Letters}, vol. 122, p. 040504, 2019.

\bibitem{liu2021}
Y. Liu, S. Arunachalam, and K. Temme, 
``A rigorous and robust quantum speed-up in supervised 
machine learning,'' \textit{Nature Physics}, vol. 17, 
pp. 1013--1017, 2021.

\bibitem{cerezo2021}
M. Cerezo et al., ``Variational quantum algorithms,'' 
\textit{Nature Reviews Physics}, vol. 3, pp. 625--644, 2021.

\bibitem{gilfuster2026}
E. Gil-Fuster et al., ``Optimal algorithmic complexity of 
inference in quantum kernel methods,'' 
arXiv:2604.15214, 2026.

\bibitem{kakavand2026}
S. Kakavand, C. Strohmeyer, and M. Schlotter, 
``Benchmarking quantum kernel SVMs against classical 
baselines on tabular data,'' arXiv:2604.18837, 2026.

\bibitem{orus2019}
R. Or\'{u}s, S. Mugel, and E. Lizaso, 
``Quantum computing for finance: Overview and prospects,'' 
\textit{Reviews in Physics}, vol. 4, p. 100028, 2019.

\bibitem{egger2020}
D. J. Egger et al., ``Quantum computing for finance,'' 
\textit{IEEE Transactions on Quantum Engineering}, 
vol. 1, pp. 1--24, 2020.

\bibitem{roy2026}
S. Roy and S. B. Alam, ``Q-SINDy: Quantum-kernel sparse 
identification of nonlinear dynamics,'' 
arXiv:2604.16779, 2026.

\bibitem{huang2026}
H.-W. Huang et al., ``Noise-enhanced quantum kernels 
on analog quantum computers,'' arXiv:2604.12476, 2026.

\bibitem{chawla2002smote}
N. V. Chawla et al., ``SMOTE: Synthetic minority 
over-sampling technique,'' 
\textit{Journal of Artificial Intelligence Research}, 
vol. 16, pp. 321--357, 2002.

\bibitem{cristianini2001}
N. Cristianini et al., ``On kernel-target alignment,'' 
in \textit{Advances in NeurIPS}, vol. 14, 2001, 
pp. 367--373.

\bibitem{ieee_cis_dataset}
IEEE Computational Intelligence Society, 
``IEEE-CIS Fraud Detection Dataset,'' Kaggle, 2019.

\bibitem{pennylane2018}
V. Bergholm et al., 
``PennyLane: Automatic differentiation of hybrid 
quantum-classical computations,'' arXiv:1811.04975, 2018.

\bibitem{qiskit2023}
IBM Quantum, ``Qiskit: An open-source framework for 
quantum computing,'' 2023.

\bibitem{sklearn2011}
F. Pedregosa et al., ``Scikit-learn: Machine learning 
in Python,'' \textit{JMLR}, vol. 12, 
pp. 2825--2830, 2011.

\end{thebibliography}

\begin{IEEEbiographynophoto}{Charan Panthangi}
is a Lead Engineer in Generative AI and Product Development 
at a Fortune 500 conglomerate,Bangalore, India. He holds 
a Post Graduate Diploma in Management from IIM Indore, 
a Master of Engineering in Applied AI and Machine Learning 
from Saveetha University, and a Certification in Quantum 
Computing from IIT Madras. He has published four 
peer-reviewed papers in IEEE Xplore on applied healthcare 
AI. His research interests include quantum machine learning, 
large language model evaluation, and applied AI for 
financial systems.
\end{IEEEbiographynophoto}

\end{document}